\section{Project Overview}
\label{sec:project-overview}

A finite dielectric slab introduces two impedance discontinuities and a propagation
phase delay. At normal incidence, the total reflected wave is the coherent sum of the
reflection at the first interface and the multiple reflections that return from the
second interface. The result depends on both relative permittivity and electrical
thickness \cite{balanis2012advanced,stutzman2013antenna}.

This study uses a fixed physical slab thickness of approximately one free-space
wavelength at 1~GHz and sweeps the dielectric relative permittivity. The analytical
model predicts a repeated set of matched conditions because the wavelength inside the
dielectric changes as $1/\sqrt{\varepsilon_r}$. A CST frequency-domain model is used
to test the predicted locations of the reflection minima.

\begin{table}[H]
\centering
\caption{Principal analytical and simulated results.}
\label{tab:principal-results}
\begin{tabularx}{\textwidth}{@{}>{\raggedright\arraybackslash}p{0.36\textwidth}X@{}}
\toprule
\textbf{Quantity} & \textbf{Result} \\
\midrule
Operating frequency & 1.000~GHz \\
Free-space wavelength & 299.792~mm analytically; 299.79~mm used in CST geometry \\
Slab thickness & $d\approx\lambda_0=299.79$~mm \\
Permittivity sweep & $\varepsilon_r=1$ to 10 in 0.25 increments, 37 points \\
Predicted matched values & $\varepsilon_r=1$, 2.25, 4, 6.25, 9 \\
CST minimum locations & Same five sweep values \\
Representative CST minimum & $\varepsilon_r=4$: $|S_{11}|=0.0104206$ or $-39.64$~dB \\
Field-validation case & $\varepsilon_r=4$ at 1~GHz \\
Evidence status & Analytical and simulated; no measured validation \\
\bottomrule
\end{tabularx}
\end{table}
